\Needspace{8\baselineskip}
\Article{177}
\textsc{Problem.} \emph{Given a sequence of forms, each a right neighbor of its predecessor, find a proper transformation of the first into any form in the sequence.}

\SourceLeaf{158}
\emph{Solution.} Let the forms be $(a,b,a')=F$; $(a',b',a'')=F'$; $(a'',b'',a''')=F''$; $(a''',b''',a'''')=F'''$ etc. Denote $\frac{b+b'}{a'}$, $\frac{b'+b''}{a''}$, $\frac{b''+b'''}{a'''}$ etc. by $h',h'',h'''$ etc., respectively. Let the indeterminates of the forms $F,F',F''$ etc. be $x,y$; $x',y'$; $x'',y''$ etc. Suppose $F$ changes
\[
\begin{array}{ll@{\qquad}l@{\qquad}l}
\text{into }F'&\text{for}&x=\alpha'x'+\beta'y',&y=\gamma'x'+\delta'y'\\
F''&\ldots&x=\alpha''x''+\beta''y'',&y=\gamma''x''+\delta''y''\\
F'''&\ldots&x=\alpha'''x'''+\beta'''y''',&y=\gamma'''x'''+\delta'''y'''\\
&&\text{etc.}&
\end{array}
\]
{\small\setlength{\arraycolsep}{3pt}
\[
\begin{array}{l@{\enspace}l@{\enspace}l}
\text{Then since }F\text{ changes into }F'\text{ for}&x=-y',&y=x'+h'y'\\
\hfill F'\text{ into }F''\text{ for}&x'=-y'',&y'=x''+h''y''\\
\hfill F''\text{ into }F'''\text{ for}&x''=-y''',&y''=x'''+h'''y'''\text{ etc. (Article\ 160)}
\end{array}
\]
}
we easily obtain the following algorithm (Article 159):
\[
\begin{array}{ll@{\qquad}ll}
\alpha'=0&\beta'=-1&\gamma'=1&\delta'=h'\\
\alpha''=\beta'&\beta''=h''\beta'-\alpha'&\gamma''=\delta'&\delta''=h''\delta'-\gamma'\\
\alpha'''=\beta''&\beta'''=h'''\beta''-\alpha''&\gamma'''=\delta''&\delta'''=h'''\delta''-\gamma''\\
\alpha''''=\beta'''&\beta''''=h''''\beta'''-\alpha'''&\gamma''''=\delta'''&\delta''''=h''''\delta'''-\gamma'''\\
&\text{etc.}&&
\end{array}
\]
or
\[
\begin{array}{ll@{\qquad}ll}
\alpha'=0&\beta'=-1&\gamma'=1&\delta'=h'\\
\alpha''=\beta'&\beta''=h''\beta'&\gamma''=\delta'&\delta''=h''\delta'-1\\
\alpha'''=\beta''&\beta'''=h'''\beta''-\beta'&\gamma'''=\delta''&\delta'''=h'''\delta''-\delta'\\
\alpha''''=\beta'''&\beta''''=h''''\beta'''-\beta''&\gamma''''=\delta'''&\delta''''=h''''\delta'''-\delta''\\
&\text{etc.}&&
\end{array}
\]

That all these transformations are proper follows easily both from their construction and from Article 159.

This very simple algorithm is convenient for calculation and is analogous to the algorithm given in Article 27; it can also be reduced to that algorithm\HistoricalNote{*)}{In the notation of Article 27, we have $\beta^n=\pm[-h'',h''',-h''''\ldots\pm h^n]$, where the alternative signs must be $-\,-$; $-\,+$; $+\,-$; $+\,+$, according as $n$ has the form $4k+0$; $1$; $2$; $3$, and
\[
\delta^n=\pm[h',-h'',h'''\ldotp\ldotp\ldotp\ldotp\ldotp\pm h^n]
\]
where the alternative signs must be $+\,-$; $+\,+$; $-\,-$; $-\,+$, according as $n$ has the form $4k+0$; $1$; $2$; $3$. But to keep this brief, we cannot develop this further; anyone can easily confirm it for themselves.}. This solution
\SourceLeaf{159}
is not restricted to forms with negative determinant: it applies to all cases, provided none of the numbers $a',a'',a'''$ etc. is $=0$.

\Needspace{8\baselineskip}
\Article{178}
\textsc{Problem.} \emph{Given two properly equivalent forms $F$, $f$ with the same negative determinant, find a proper transformation of one into the other.}

\emph{Solution.} Suppose the form $F$ is $(A,B,A')$, and that the method of Article 171 has given the sequence of forms $(A',B',A'')$, $(A'',B'',A''')$ etc. up to $(A^m,B^m,A^{m+1})$, which is reduced. Similarly, suppose $f$ is $(a,b,a')$, and the same method has given the sequence $(a',b',a'')$, $(a'',b'',a''')$ up to $(a^n,b^n,a^{n+1})$, which is reduced. There are then two possible cases.

I. The forms $(A^m,B^m,A^{m+1})$, $(a^n,b^n,a^{n+1})$ are either identical, or opposite and also ambiguous. Then the forms $(A^{m-1},B^{m-1},A^m)$, $(a^n,-b^{n-1},a^{n-1})$ are neighbors (where $A^{m-1}$ denotes the next-to-last term of the sequence $A,A',A''\ldots A^m$, and $B^{m-1}$, $a^{n-1}$, $b^{n-1}$ have similar meanings). For $A^m=a^n$, $B^{m-1}\equiv-B^m\ (\mathrm{mod.}\ A^m)$, $b^{n-1}\equiv-b^n\ (\mathrm{mod.}\ a^n\text{ or }A^m)$, giving $B^{m-1}-b^{n-1}\equiv b^n-B^m$ and so $\equiv0$ if the forms $(A^m,B^m,A^{m+1})$, $(a^n,b^n,a^{n+1})$ are identical, and $\equiv2b^n$, so $\equiv0$, if they are opposite and ambiguous. Thus, in the sequence of forms
\[
\begin{gathered}
(A,B,A'),\ (A',B',A'')\ldotp\ldotp\ldotp\ldotp(A^{m-1},B^{m-1},A^m),\\
(a^n,-b^{n-1},a^{n-1}),\ (a^{n-1},-b^{n-2},a^{n-2})\ldots(a',-b,a),\ (a,b,a')
\end{gathered}
\]
each form is a neighbor of its predecessor. By the preceding article, we can therefore find a proper transformation of the first, $F$, into the last, $f$.

II. The forms $(A^m,B^m,A^{m+1})$, $(a^n,b^n,a^{n+1})$ are not identical but opposite, and also $A^m=A^{m+1}=a^n=a^{n+1}$. Then the sequence of forms
\[
\begin{gathered}
(A,B,A'),\ (A',B',A'')\ldotp\ldotp\ldotp\ldotp(A^m,B^m,A^{m+1}),\\
(a^n,-b^{n-1},a^{n-1}),\ (a^{n-1},-b^{n-2},a^{n-2})\ldotp\ldotp\ldotp\ldotp(a',-b,a),\ (a,b,a')
\end{gathered}
\]
has the same property. For $A^{m+1}=a^n$, and $B^m-b^{n-1}=-(b^n+b^{n-1})$ is divisible by $a^n$. So the preceding article gives a proper transformation of the first form $F$ into the last $f$.

\SourceLeaf{160}
\emph{Example.} For the forms $(23.\,38,63)$, $(15,20,27)$, we have the sequence
\[
\begin{gathered}
(23,38,63).\ (63,25,10),\ (10,5,3),\ (3,1.\,2),\\
(2.\,-7,27),\ (27,-20,15),\ (15,20,27)
\end{gathered}
\]
so
\[
h'=1,\quad h''=3,\quad h'''=2,\quad h''''=-3,\quad h'''''=-1,\quad h''''''=0.
\]
This gives the following transformation of the form $23xx+76xy+63yy$ into $15tt+40tu+27uu$: $x=-13t-18u$, $y=8t+11u$.

This solution easily gives the solution of the problem: \emph{If the forms $F,f$ are improperly equivalent, find an improper transformation of the form $F$ into $f$.} For let $f=att+2btu+a'uu$; then the opposite form $app-2bpq+a'qq$ is properly equivalent to the form $F$. Find a proper transformation of the form $F$ into that form, $x=\alpha p+\beta q$, $y=\gamma p+\delta q$. Clearly $F$ changes into $f$ on putting $x=\alpha t-\beta u$, $y=\gamma t-\delta u$, and this transformation is improper.

So if the forms $F,f$ are both properly and improperly equivalent, we can find both a proper and an improper transformation.

\Needspace{8\baselineskip}
\Article{179}
\textsc{Problem.} \emph{If the forms $F$, $f$ are equivalent, find all transformations of the form $F$ into $f$.}

\emph{Solution.} If the forms $F$, $f$ are equivalent in only one way, \emph{that is}, only properly or only improperly, use the preceding article to find one transformation of the form $F$ into $f$. Clearly there can be no other transformations except those similar to this one. If the forms $F$, $f$ are both properly and improperly equivalent, find two transformations, one proper and the other improper. Now take the form $F=(A,B,C)$, with $BB-AC=-D$, and let the greatest common divisor of the numbers $A,2B,C$ be $=m$. Article 162 shows that, in the first case, all transformations of the form $F$ into $f$ can be derived from one transformation; in the second, all proper transformations can be derived from the proper one and all improper transformations from the improper one, provided we have all solutions of the equation $tt+Duu=mm$. Once these have been found, the problem is solved.

But we have $D=AC-BB$, $4D=4AC-4BB$, so $\frac{4D}{mm}=4\frac{AC}{mm}-(\frac{2B}{m})^2$ is an integer. Now if

1) $\frac{4D}{mm}>4$, we have $D>mm$. So in $tt+Duu=mm$, $u$ must be $=0$, and $t$ can have no values other than $+m$ and $-m$. Thus,
\SourceLeaf{161}
if $F$, $f$ are equivalent in only one way and one transformation is
\[
x=\alpha x'+\beta y',\quad y=\gamma x'+\delta y'
\]
there can be no others except this one, which comes from $t=m$ (Article 162), and
\[
x=-\alpha x'-\beta y',\quad y=-\gamma x'-\delta y'
\]
If $F$, $f$ are both properly and improperly equivalent, and we have a proper transformation
\[
x=\alpha x'+\beta y',\quad y=\gamma x'+\delta y'
\]
and an improper one
\[
x=\alpha'x'+\beta'y',\quad y=\gamma'x'+\delta'y'
\]
there is no other proper transformation except the former (from $t=m$) and
\[
x=-\alpha x'-\beta y',\ y=-\gamma x'-\delta y'\ (\text{from }t=-m)
\]
and likewise no improper transformation except
\[
\begin{gathered}
x=\alpha'x'+\beta'y',\quad y=\gamma'x'+\delta'y';\\
\text{and }x=-\alpha'x'-\beta'y',\quad y'=-\gamma'x'-\delta'y'
\end{gathered}
\]

2) If $\frac{4D}{mm}=4$, or $D=mm$, the equation $tt+Duu=mm$ has four solutions $t,u;=m,0;\ -m,0;\ 0,1;\ 0,-1$. So if $F$, $f$ are equivalent in only one way and one transformation is
\[
x=\alpha x'+\beta y',\quad y=\gamma x'+\delta y'
\]
there are exactly \emph{four} transformations,
\[
x=\pm\alpha x'\pm\beta y',\quad y=\pm\gamma x'\pm\delta y'
\]
\[
\begin{aligned}
x&=\mp\frac{\alpha B+\gamma C}{m}x'\mp\frac{\beta B+\delta C}{m}y',\\
y&=\pm\frac{\alpha A+\gamma B}{m}x'\pm\frac{\beta A+\delta B}{m}y'
\end{aligned}
\]

If $F$, $f$ are equivalent in two ways, or there is another transformation dissimilar to the given one, this also gives four transformations dissimilar to the former four, making \emph{eight} transformations. --- It is also easy to prove that in this case $F$, $f$ really are always equivalent in both ways. For since $D=mm=AC-BB$, $m$ divides $B$ too. The form $(\frac{A}{m},\frac{B}{m},\frac{C}{m})$ has determinant $=-1$, so it is equivalent to the form $(1,0,1)$ or to $(-1,0,-1)$.
\SourceLeaf{162}
It is easy to see that the same transformation that changes $(\frac{A}{m},\frac{B}{m},\frac{C}{m})$ into $(\pm1,0,\pm1)$ changes the form $(A,B,C)$ into $(\pm m,0,\pm m)$, which is ambiguous. So the form $(A,B,C)$, being equivalent to an ambiguous form, is both properly and improperly equivalent to every form equivalent to it.

3) If $\frac{4D}{mm}=3$, or $4D=3mm$, $m$ is even, and all solutions of the equation $tt+Duu=mm$ are these six:
\[
t,u;=m,0;\;-m,0;\;\tfrac12m,1;\;-\tfrac12m,-1;\;\tfrac12m,-1;\;-\tfrac12m,1
\]
So if we have two dissimilar transformations of the form $F$ into $f$,
\[
\begin{aligned}
x&=\alpha x'+\beta y',& y&=\gamma x'+\delta y'\\
x&=\alpha'x'+\beta'y',& y&=\gamma'x'+\delta'y'
\end{aligned}
\]
we have twelve transformations, namely six similar to the first,
\[
x=\pm\alpha x'\pm\beta y',\quad y=\pm\gamma x'\pm\delta y'
\]
\[
\begin{aligned}
x&=\pm(\tfrac12\alpha-\tfrac{\alpha B+\gamma C}{m})x'\pm(\tfrac12\beta-\tfrac{\beta B+\delta C}{m})y'\\
y&=\pm(\tfrac12\gamma+\tfrac{\alpha A+\gamma B}{m})x'\pm(\tfrac12\delta+\tfrac{\beta A+\delta B}{m})y'
\end{aligned}
\]
\[
\begin{aligned}
x&=\pm(\tfrac12\alpha+\tfrac{\alpha B+\gamma C}{m})x'\pm(\tfrac12\beta+\tfrac{\beta B+\delta C}{m})y'\\
y&=\pm(\tfrac12\gamma-\tfrac{\alpha A+\gamma B}{m})x'\pm(\tfrac12\delta-\tfrac{\beta A+\delta B}{m})y'
\end{aligned}
\]
and six similar to the second, obtained from these by replacing $\alpha$, $\beta$, $\gamma$, $\delta$ with $\alpha'$, $\beta'$, $\gamma'$, $\delta'$.

That in this case $F,f$ are always equivalent in both ways is proved as follows. The form $(\frac{2A}{m},\frac{2B}{m},\frac{2C}{m})$ has determinant $=-\frac{4D}{mm}=-3$, so this form (Article 176) is equivalent either to the form $(\pm1,0,\pm3)$ or to $(\pm2,\pm1,\pm2)$. From this it is easy to see that the form $(A,B,C)$ is equivalent either to the form $(\pm\tfrac12m,0,\pm\tfrac32m)$ or to $(\pm m,\tfrac12m,\pm m)$\HistoricalNote{*)}{It can be proved that the form $(A,B,C)$ is necessarily equivalent to the latter, but this is not needed here.}, both of which are ambiguous, and so is equivalent in both ways to every equivalent form.

4) If we assume $\frac{4D}{mm}=2$, we get $(\frac{2B}{m})^2=4\frac{AC}{mm}-2$, and so $\equiv2\ (\mathrm{mod.}\ 4)$.
\SourceLeaf{163}
But no square can be $\equiv2\ (\mathrm{mod.}\ 4)$, so this case cannot occur.

5) Assuming $\frac{4D}{mm}=1$, we get $(\frac{2B}{m})^2=4\frac{AC}{mm}-1\equiv-1\ (\mathrm{mod.}\ 4)$. Since this is impossible, this case cannot occur either.

Since $D$ is neither $=0$ nor negative, no cases other than those listed are possible.

\Needspace{8\baselineskip}
\Article{180}
\textsc{Problem.} \textit{Find all representations of a given number $M$ by the form $axx+2bxy+cyy\ldots F$ with negative determinant $-D$, in which $x,y$ take relatively prime values.}

\textit{Solution.} Article 154 shows that $M$ cannot be represented in the required way unless $-D$ is a quadratic residue of $M$. First find all distinct (\textit{that is}, incongruent) values of the expression $\surd-D\ (\mathrm{mod.}\ M)$, say $N,-N,N',-N',N'',-N''$ etc. To make the calculation simpler, all $N,N'$ etc. can be chosen so that they are not $>\tfrac12M$. Since each representation must belong to one of these values, consider each value separately.

If the forms $F,(M,N,\frac{D+NN}{M})$ are not properly equivalent, there can be no representation of $M$ belonging to the value $N$ (Article 168). If they are, find a proper transformation of the form $F$ into
\[
Mx'x'+2Nx'y'+\frac{D+NN}{M}y'y'
\]
say
\[
x=\alpha x'+\beta y',\quad y=\gamma x'+\delta y'
\]
Then $x=\alpha,y=\gamma$ is a representation of the number $M$ by $F$ belonging to $N$. Let the greatest common divisor of the numbers $A,2B,C,=m$. We distinguish three cases (preceding article):

1) If $\frac{4D}{mm}>4$, there are no representations belonging to $N$ other than these \textit{two}: $x=\alpha,y=\gamma$; $x=-\alpha,y=-\gamma$ (Articles 169, 179).

2) If $\frac{4D}{mm}=4$, there are \textit{four} representations
\[
x=\pm\alpha,\quad y=\pm\gamma;\quad x=\mp\frac{\alpha B+\gamma C}{m},\quad y=\pm\frac{\alpha A+\gamma B}{m}
\]
3) If $\frac{4D}{mm}=3$, there are \textit{six} representations
\SourceLeaf{164}
\[
\begin{aligned}
x&=\pm\alpha,&y&=\pm\gamma;\\
x&=\pm(\tfrac12\alpha-\tfrac{\alpha B+\gamma C}{m}),&y&=\pm(\tfrac12\gamma+\tfrac{\alpha A+\gamma B}{m})\\
x&=\pm(\tfrac12\alpha+\tfrac{\alpha B+\gamma C}{m}),&y&=\pm(\tfrac12\gamma-\tfrac{\alpha A+\gamma B}{m})
\end{aligned}
\]
Representations belonging to the values $-N.N',-N'$ etc. must be found in the same way.

\Needspace{25\baselineskip}
\Article{181}
Finding the representations of the number $M$ by the form $F$ in which $x,y$ take values that are not relatively prime can easily be reduced to the case already considered. Suppose such a representation is obtained on putting $x=\mu e,y=\mu f$, where $\mu$ is the greatest common divisor of $\mu e,\mu f$, that is, $e,f$ are relatively prime. Then $M=\mu\mu(Aee+2Bef+Cff)$, which is therefore divisible by $\mu\mu$. The substitution $x=e,y=f$ is a representation of the number $\frac{M}{\mu\mu}$ by the form $F$ in which $x,y$ take relatively prime values. So if $M$ is divisible by no square (other than 1), \textit{for example}, if it is prime, there are no such representations of $M$. If $M$ has square divisors, let them be $\mu\mu,\nu\nu,\pi\pi$ etc. First find all representations of the number $\frac{M}{\mu\mu}$ by the form $(A,B,C)$ in which $x,y$ take relatively prime values. Multiplying these values by $\mu$ gives all representations of $M$ in which the greatest common divisor of the numbers $x,y$ is $\mu$. Similarly, all representations of $\frac{M}{\nu\nu}$ in which the values of $x,y$ are relatively prime give all representations of $M$ in which the greatest common divisor of the values of $x,y$ is $\nu$, etc.

So the preceding rules find all representations of a given number by a given form with negative determinant.
